\subsection{Receiver-sufficient prefixes}\label{subsec:receiver-sufficient-prefix}

The family distinction criterion has a useful architectural specialization. Suppose the whole reachable downstream state is determined by an earlier-cut interface $P$, while every downstream receiver interface retains a local coordinate whose earlier-cut distinction shadow is exactly that of $P$. Then no downstream parameterization can create a finer predecessor distinction relative to the earlier cut.

\begin{definition}[Receiver-sufficient prefix]\label{def:receiver-sufficient-prefix-v13}
Let
\[
P:\Omega\twoheadrightarrow U,
\qquad
A:\Omega\to X,
\qquad
L_j:X\to V_j.
\]
We call $P$ receiver-sufficient for $A$ at receiver $j$ when:
\begin{enumerate}[label=(\roman*)]
\item the whole reachable downstream state factors through $P$, so there exists $\widetilde A:U\to X$ with
\[
A=\widetilde A P;
\]
\item the recut local receiver state has exactly the $P$ distinction shadow,
\[
L_jA\equiv_DP.
\]
\end{enumerate}
\end{definition}

Consider any parameterized downstream receiver family whose actual presented interfaces retain this local coordinate,
\[
Q_{j,\theta}(x)=\bigl(L_j(x),E_{j,\theta}(x)\bigr)
\]
on their realized images.

\begin{proposition}[Receiver-sufficient distinction collapse]\label{prop:receiver-sufficient-collapse-v13}
If $P$ is receiver-sufficient for $A$ at receiver $j$, then for every parameter state $\theta$,
\[
\boxed{Q_{j,\theta}A\equiv_DP.}
\]
Consequently,
\[
\boxed{
\mathcal C_A(\mathcal Q_j)=\{[P]_D\},
\qquad
J_{A,\mathcal Q_j}\equiv_DP.
}
\]
Moreover, for each $\theta$ there is a unique bijection on realized images
\[
T_\theta:U\xrightarrow{\cong}\im(Q_{j,\theta}A)
\]
such that
\[
\boxed{Q_{j,\theta}A=T_\theta P.}
\]
\end{proposition}

\begin{proof}
Because $A=\widetilde A P$, every composite $Q_{j,\theta}A$ factors through $P$, so $Q_{j,\theta}A\preceq_DP$. Conversely, equality under $Q_{j,\theta}A$ implies equality of its retained local coordinate $L_jA$. Since $L_jA\equiv_DP$, equality under the receiver interface implies equality under $P$, so $P\preceq_DQ_{j,\theta}A$. The kernels are equal. The class-set and envelope conclusions follow, and the unique $T_\theta$ follows from \Cref{thm:unique-conversion-v13}.
\end{proof}

\begin{corollary}[Injective receiver-sufficient collapse]\label{cor:injective-receiver-sufficient-v14}
Under the hypotheses of \Cref{prop:receiver-sufficient-collapse-v13}, suppose additionally that $P$ is injective. Then:
\begin{enumerate}[label=(\roman*)]
\item $A$ is injective;
\item every effective receiver interface $Q_{j,\theta}A$ is injective;
\item every effective distinction shadow is the diagonal relation
\[
\Delta_\Omega:=\{(\omega,\omega):\omega\in\Omega\};
\]
\item any distinction-class collapse supplied by the proposition cannot be attributed to loss of predecessor distinctions by the prefix $A$.
\end{enumerate}
\end{corollary}

\begin{proof}
Since $L_jA\equiv_DP$ and $P$ is injective,
\[
\ker(L_jA)=\ker P=\Delta_\Omega,
\]
so $L_jA$ is injective. If $A(\omega)=A(\omega')$, then $L_jA(\omega)=L_jA(\omega')$, hence $\omega=\omega'$. Therefore $A$ is injective. The receiver-sufficient collapse proposition gives
\[
\ker(Q_{j,\theta}A)=\ker P=\Delta_\Omega
\]
for every $\theta$, proving the remaining claims.
\end{proof}

The distinction between the two results matters. Receiver sufficiency alone can arise because an upstream prefix discarded predecessor information. The injective corollary isolates a different phenomenon: the prefix can preserve every predecessor distinction while changing which downstream schema distinctions remain visible at the chosen receiver-access grain. Section~\ref{subsec:attention-context-witness} gives an exact attention instance.
